% !PRINTED 46
\weberpage{46}
%<<CONT
\weq{11}{\begin{vmatrix} 1, & x_1, & x_1^2, & \sqrt{f(x_1)} \\ 1, & x_2, & x_2^2, & \sqrt{f(x_2)} \\ 1, & x_3, & x_3^2, & \sqrt{f(x_3)} \\ 1, & x_4, & x_4^2, & \sqrt{f(x_4)} \end{vmatrix} = 0.}

From this equation one can express $x_1$ and $\sqrt{f(x_1)}$ \gesperrt{rationally}
in terms of $x_2, \sqrt{f(x_2)}$; $x_3, \sqrt{f(x_3)}$; $x_4, \sqrt{f(x_4)}$; for
$x_1$ is the root of a biquadratic equation whose three
other roots are $x_2, x_3, x_4$, and whose coefficients depend
rationally on $x_2, \sqrt{f(x_2)}$; $x_3, \sqrt{f(x_3)}$; $x_4, \sqrt{f(x_4)}$, and $\sqrt{f(x_1)}$ can
be represented rationally in terms of $x_1$ by means of (11). If one sets $x_4$ equal to an
arbitrary constant, the differential equation
(10) passes into (1), and (11) yields its integral.

We apply this theorem first to the normal form % German source: "dies Theorem"
and set
\[ f(x) = x(1-x)(1-\varkappa^2 x). \]
If we assume $x_4 = 0$ in (11), then this
equation reduces, on division by $\sqrt{x_1 x_2 x_3}$, to:
\weq{12}{\begin{vmatrix} \sqrt{x_1}, & x_1\sqrt{x_1}, & \sqrt{(1-x_1)(1-\varkappa^2 x_1)} \\ \sqrt{x_2}, & x_2\sqrt{x_2}, & \sqrt{(1-x_2)(1-\varkappa^2 x_2)} \\ \sqrt{x_3}, & x_3\sqrt{x_3}, & \sqrt{(1-x_3)(1-\varkappa^2 x_3)} \end{vmatrix} = 0,}
or, if we develop the determinant according to the elements of the first
row:
\weq{13}{\sqrt{x_1}\,(a + bx_1) + c\sqrt{(1-x_1)(1-\varkappa^2 x_1)} = 0,}
in which, for abbreviation, we have set:
\weq{14}{\begin{aligned}
a &= x_2\sqrt{x_2}\sqrt{(1-x_3)(1-\varkappa^2 x_3)} - x_3\sqrt{x_3}\sqrt{(1-x_2)(1-\varkappa^2 x_2)}, \\
b &= -\sqrt{x_2}\sqrt{(1-x_3)(1-\varkappa^2 x_3)} + \sqrt{x_3}\sqrt{(1-x_2)(1-\varkappa^2 x_2)}, \\
c &= -(x_2-x_3)\sqrt{x_2}\sqrt{x_3}.
\end{aligned}}

If we make (13) rational, we obtain the cubic
equation:
\[ x(a+bx)^2 - c^2(1-x)(1-\varkappa^2 x) = 0, \]
whose three roots are $x_1, x_2, x_3$. Accordingly we have identically
\weq{15}{b^2(x-x_1)(x-x_2)(x-x_3) = x(a+bx)^2 - c^2(1-x)(1-\varkappa^2 x),}
%CONT>>
